\documentclass[10pt,a4paper]{amsart}
\usepackage[margin=19mm]{geometry}
\usepackage{amsmath,amssymb,mathtools}
\usepackage[scr=rsfs,cal=euler]{mathalfa}
\usepackage{xfrac}
\usepackage[unicode,hidelinks,bookmarks=false]{hyperref}
\newcommand{\ie}{i.e.\ }
\newcommand{\cf}{cf.\ }
\newcommand{\resp}{respectively\ }
\newcommand{\Subsection}{\subsection}
\providecommand{\nobreakdash}{\nobreak\hbox{-}}
\setlength{\emergencystretch}{2em}
\multlinegap0pt
\usepackage{bm}
\usepackage{bbm}
\usepackage{dsfont}
\usepackage{xspace}
\usepackage{cancel}
\usepackage{mathtools}
\usepackage{amsthm}
\makeatletter
\@ifundefined{theorem}{\newtheorem{theorem}{Theorem}}{}
\@ifundefined{proposition}{\newtheorem{proposition}[theorem]{Proposition}}{}
\makeatother
\makeatletter
\def\H{{\mathcal H}}
\expandafter\ifx\csname Proof\endcsname\relax
\newenvironment{Proof}
% {\noindent{\it Proof\/}}{{ \hfill $\Box$}\smallskip\par}

\fi
\expandafter\ifx\csname dedication\endcsname\relax
\newenvironment{dedication}
  {\clearpage           % we want a new page
   \thispagestyle{empty}% no header and footer
   \vspace*{\stretch{1}}% some space at the top 
  \section*{Dedication} % <==========================================
   \itshape             % the text is in italics
   \raggedleft          % flush to the right margin
  }
  {\par % end the paragraph
   \vspace{\stretch{3}} % space at bottom is three times that at the top
   \clearpage           % finish off the page
  }
\fi
\expandafter\ifx\csname definition\endcsname\relax
\newenvironment{definition}
% {\smallskip\noindent{\bf Definition\/}:}{\smallskip\par}
% \newenvironment{statement}
% {\smallskip\noindent{\bf Statement\/}.}{\smallskip\par}
% \newenvironment{example}
% {\smallskip\noindent{\bf Example\/}.}{\smallskip\par}
% \newenvironment{remark}
% {\smallskip\noindent{\bf Remark\/}.}{\smallskip\par}
% \newenvironment{remarks}
% {\smallskip\noindent{\bf Remarks\/}.}{\smallskip\par}
% \newenvironment{proof}
% {\noindent{\it Proof\/}.}{{ \hfill $\Box$}\smallskip\par}
% \newenvironment{Proof}
% {\noindent{\it Proof\/}}{{ \hfill $\Box$}\smallskip\par}

\fi
\expandafter\ifx\csname example\endcsname\relax
\newenvironment{example}
% {\smallskip\noindent{\bf Example\/}.}{\smallskip\par}
% \newenvironment{remark}
% {\smallskip\noindent{\bf Remark\/}.}{\smallskip\par}
% \newenvironment{remarks}
% {\smallskip\noindent{\bf Remarks\/}.}{\smallskip\par}
% \newenvironment{proof}
% {\noindent{\it Proof\/}.}{{ \hfill $\Box$}\smallskip\par}
% \newenvironment{Proof}
% {\noindent{\it Proof\/}}{{ \hfill $\Box$}\smallskip\par}

\fi
\newcommand{\id}{\mathrm{id}}
\expandafter\ifx\csname namelist\endcsname\relax
\newenvironment{namelist}[1]{%
\begin{list}{}
{\let\makelabel\namelistlabel
\settowidth{\labelwidth}{#1}
\setlength{\leftmargin}{1.1\labelwidth}}
}{%
\end{list}}
\fi
\def\p{\partial }
\expandafter\ifx\csname proof\endcsname\relax
\newenvironment{proof}
% {\noindent{\it Proof\/}.}{{ \hfill $\Box$}\smallskip\par}
% \newenvironment{Proof}
% {\noindent{\it Proof\/}}{{ \hfill $\Box$}\smallskip\par}

\fi
\expandafter\ifx\csname remark\endcsname\relax
\newenvironment{remark}
% {\smallskip\noindent{\bf Remark\/}.}{\smallskip\par}
% \newenvironment{remarks}
% {\smallskip\noindent{\bf Remarks\/}.}{\smallskip\par}
% \newenvironment{proof}
% {\noindent{\it Proof\/}.}{{ \hfill $\Box$}\smallskip\par}
% \newenvironment{Proof}
% {\noindent{\it Proof\/}}{{ \hfill $\Box$}\smallskip\par}

\fi
\expandafter\ifx\csname remarks\endcsname\relax
\newenvironment{remarks}
% {\smallskip\noindent{\bf Remarks\/}.}{\smallskip\par}
% \newenvironment{proof}
% {\noindent{\it Proof\/}.}{{ \hfill $\Box$}\smallskip\par}
% \newenvironment{Proof}
% {\noindent{\it Proof\/}}{{ \hfill $\Box$}\smallskip\par}

\fi
\expandafter\ifx\csname statement\endcsname\relax
\newenvironment{statement}
% {\smallskip\noindent{\bf Statement\/}.}{\smallskip\par}
% \newenvironment{example}
% {\smallskip\noindent{\bf Example\/}.}{\smallskip\par}
% \newenvironment{remark}
% {\smallskip\noindent{\bf Remark\/}.}{\smallskip\par}
% \newenvironment{remarks}
% {\smallskip\noindent{\bf Remarks\/}.}{\smallskip\par}
% \newenvironment{proof}
% {\noindent{\it Proof\/}.}{{ \hfill $\Box$}\smallskip\par}
% \newenvironment{Proof}
% {\noindent{\it Proof\/}}{{ \hfill $\Box$}\smallskip\par}

\fi
\makeatother
\title[Orbifold Saito theory of A and D type singularities]{Orbifold Saito theory of A and D type singularities}
\author{Alexey Basalaev, Anton Rarovskii}
\date{Source: arXiv:2507.07609v1 (CC BY 4.0)}
\begin{document}
\maketitle

\noindent\emph{Source and licence.} Orbifold Saito theory of A and D type singularities. Version arXiv:2507.07609v1 (CC BY 4.0), distributed under the Creative Commons Attribution 4.0 International licence, as recorded in the arXiv licence metadata for this exact version and shown on the abstract page https://arxiv.org/abs/2507.07609v1. Full rights notice: licenses/arxiv-2507.07609-CC-BY-4.0.txt.\par\smallskip

\noindent\textit{Editorial scope.} Counted unit: the Proposition whose source label is \texttt{EF4} in the article (source0.tex lines 1688--1724, source label \texttt{EF4}), delivered together with the article's own complete original proof of it (source0.tex lines 1726--1821, one \texttt{proof} environment of 96 lines). No mathematics is written, repaired or re-derived: statement and proof are the article's own text line for line. Required context retained uncounted: none. Every definition, display and label of those lines is retained unchanged. Omitted from this package and not redistributed: 1--1687, 1725--1725, 1822--1934. Nothing inside a retained excerpt is omitted or altered: every retained body is delivered exactly as the article's lines are written, so no omission receipt of any kind accompanies this package. Citations: the retained excerpts carry no citation command at all, so no bibliography entry of the article is transcribed and no reference is redistributed. Reference adaptation: none was needed and none was made; every reference inside the delivered unit resolves inside it, so no unresolved reference or citation is produced. Printed slips kept literal: no printed word, symbol, hypothesis or number of the counted statement or of its proof was altered. Layout and class: the article's own class is \texttt{amsart} with packages including hyperref, nameref, zref-xr, amsmath, amsthm, amssymb, amsfonts, amscd, epsfig, xy, tabularx, booktabs, float, mathtext; the wrapper is the lane's pinned \texttt{amsart} editorial wrapper at 10pt on A4, which restates the theorem-like environments the retained text needs and adds the article's own macro definitions that the retained text needs (\texttt{\textbackslash H}, \texttt{\textbackslash id}, \texttt{\textbackslash p}), with extra wrapper packages (bm, bbm, dsfont, xspace, cancel, mathtools). The delivered numbering is the unit's own. Assets: no retained span contains a figure environment, a graphics-inclusion command, a PDF-page inclusion, a table or a TikZ picture; no image file is copied into this package, the package declares no asset dependency and no image file is redistributed with it. Licence: version arXiv:2507.07609v1 of this article is distributed under the Creative Commons Attribution 4.0 International licence, as recorded in the arXiv licence metadata for this exact version and shown on the abstract page https://arxiv.org/abs/2507.07609v1. A full rights notice is delivered at licenses/arxiv-2507.07609-CC-BY-4.0.txt. Characters: every retained excerpt is pure ASCII, so the delivered engine, XeLaTeX with its default Latin Modern text font, reports no missing character.
\begin{proposition}{\label{EF4}}
    For the pair $(f,G)$ as above, the action of orbifold Gauss-Manin connection reads
\begin{enumerate}
\item If $\alpha+\beta \le 2k$, then
\[
\nabla_{\frac{\partial}{\partial v^\alpha_{\id}}}^{(f,G)}[x_3^{2\beta}\xi_{\id}] = \nabla^{F^G}_{\frac{\partial}{\partial v^\alpha_{\id}}}[x_3^{2\beta}]\xi_{\id}.
\]
If $\alpha+\beta = 2k+1$, then
\[
\nabla_{\frac{\partial}{\partial v^\alpha_{\id}}}^{(f,G)}[x_1^{2\beta}\xi_{\id}] = \nabla^{F^G}_{\frac{\partial}{\partial v^\alpha_{\id}}}[x_1^{2\beta}]\xi_{\id} - \frac{2v_g^0}{2k+1} \left[\frac{\p F^G}{\p v_\id^{\alpha+\beta-2k-1}}\right]
\xi_{g}.
\]
If $\alpha+\beta > 2k+1$, then
\[
\nabla_{\frac{\partial}{\partial v^\alpha_{\id}}}^{(f,G)}[x_1^{2\beta}\xi_{\id}] = \nabla^{F^G}_{\frac{\partial}{\partial v^\alpha_{\id}}}[x_1^{2\beta}]\xi_{\id} + \frac{(v_g^0)^2}{2k+1} 
\left[\frac{\p F^G}{\p v_\id^{\alpha+\beta-2k-2}}\right]\xi_{\id}.
\]

\item If $1 \le \beta \le 2k$, then we have
\[
\nabla_{\frac{\partial}{\partial v^0_{g}}}^{(f,G)}[x_3^{2\beta}\xi_{\id}] = \nabla_{\frac{\partial}{\partial v^\beta_{\id}}}^{(f,G)}[\xi_{g}] = - \frac{v_g^0}{2}\left[\frac{\partial F^G}{\partial v^{\beta-1}_{\id}}\right]\xi_{\id},
\]
and
\[
\nabla_{\frac{\partial}{\partial v^0_{g}}}^{(f,G)}[1 \cdot \xi_{\id}] 
= \nabla_{\frac{\partial}{\partial v^0_{\id}}}^{(f,G)}[\xi_{g}] 
= \xi_g,
\]
\item 
\[
        \nabla_{\frac{\partial}{\partial v^0_{g}}}^{(f,G)}[\xi_{g}] = 
        c_g \nabla_{\frac{\partial}{\partial v^0_{\id}}}^{(f,G)}[H^F_g \xi_{\id}](\widetilde v_{id})
    \]
where $\widetilde v_{\id}^{\alpha} := v_{\id}^{\alpha} \frac{2k + 1}{2\alpha -1}$.    

\end{enumerate}
\end{proposition}

\begin{proof}
 (i) In this case we can not apply the map $p$ and we prove the proposition by explicit calculations via the map $\Psi_{\bar{S}}$:
\[
 \nabla_{\frac{\partial}{\partial v^\alpha_{\id}}}^{(f,G)}[x_3^{2\beta}\xi_{\id}] = \Psi_{\bar{S}}[y_3^{\alpha+\beta}]
\]
If $0 \le \alpha+\beta \le 2k$, $\Psi_{\bar{S}} [y_3^{\alpha+\beta}]_{\bar{F}} = [x_3^{2(\alpha+\beta)}]\xi_{\id} = \nabla^{F^G}_{\frac{\partial}{\partial v^\alpha_\id}}[x_1^{2\beta}]\xi_{\id}$ and the statement holds. For $2k+1\le \alpha + \beta \le 4k$ we need to take into account the relations in the Brieskorn lattice $\H_{\bar F}$.
Denote $p := \alpha+\beta - (2k+1)$. Then $0 \le p \le 2k-1$. Applying $D_{\bar{F}}$ to $-[y_3^{p+k}dy_1\wedge dy_3]$ and $[y_3^{p+1}dy_1\wedge dy_2]$ we have
\[
[2y_2y_3^{p+k+1} + y_3^{\alpha+\beta} + \bar{v}_{0} y_3^{p+k}]_{\bar{F}} = [2y_3^{p+k}(y_2 y_3+ \frac{\bar{v}_0}{2}) + y_3^{\alpha+\beta}]_{\bar{F}}= 0,
\]
\[
[z(p+1)y_3^{p} + k y_3^{p+k}(y_2y_3 + \frac{\bar{v}_0}{2}) +  y_3^{p+1}(y_2^2 + y_2y_3^k) + \sum_{\gamma=0}^{2k}\gamma\bar{s}_\gamma y_3^{\gamma +p}]_{\bar{F}} = 0.
\]
These relations give us
\[
[y_3^{\alpha+\beta}]_{\bar{F}} = \frac{2}{k}[z(p+1)y_3^p + y_3^{p+1}(y_2^2 + y_2y_3^k) + \sum_{\gamma=0}^{2k}\gamma\bar{s}_\gamma y_3^{\gamma +p}]_{\bar{F}}.
\]
We also note that
\[
[y_3^{p+1}(y_2^2 + y_2y_3^k)]_{\bar{F}} = [y_3^{p+1}((y_2 + \frac{1}{2}y_3^k)^2 - \frac{1}{4}y_3^{2k})]_{\bar{F}} = [y_3^{p+1}(y_2 + \frac{1}{2}y_3^k)^2 - \frac{1}{4}y_3^{\alpha+\beta}]_{\bar{F}}
\]
which gives us
\[
[y_3^{\alpha+\beta}]_{\bar{F}} = \frac{4}{2k+1}[z(p+1)y_3^p + y_3^{p+1}(y_2 + \frac{1}{2}y_3^k)^2 + \sum_{\gamma=0}^{2k}\gamma\bar{s}_\gamma y_3^{\gamma +p}]_{\bar{F}}.
\]
Finally, applying $D_{\bar{F}}$ to $[y_3^{p}(y_2+\frac{1}{2}y_3^k) dy_1 \wedge dy_3]$ we have
\[
[z y_3^p + 2y_3^{p+1}(y_2+\frac{1}{2}y_3^k)^2 + \bar{v}_0 y_3^{p}(y_2+\frac{1}{2}y_3^k)]_{\bar{F}} = 0.
\]
This relation gives
\[
[y_3^{\alpha+\beta}]_{\bar{F}} = \frac{4}{2k+1}[z(p+\frac{1}{2})y_3^p + \sum_{\gamma=0}^{2k}\gamma\bar{s}_\gamma y_3^{\gamma +p}]_{\bar{F}} - \frac{4}{2k+1}[\frac{\bar{v}_0}{2}y_3^{p}(y_2+\frac{1}{2}y_3^k)]_{\bar{F}}.
\]
If $p > 0$ we need to express the last summand as well. Applying $D_{\bar F}$ to $[y_3^{p-1} dy_1 \wedge dy_3]$ we have
 \[
    [ 2y_3^p(y_2 + \frac{1}{2}y_3^k) + \bar v_0 y_3^{p-1}]_{\bar F} = 0,
 \]
 what gives the expression for the second term.

To complete the proof we show that the first term is determined by $\nabla^{F^G}_{\frac{\partial}{\partial v^\alpha_{\id}}}[x_3^{2\beta}]\xi_{\id}$. Note that $\nabla^{F^G}_{\frac{\partial}{\partial v^\alpha_{\id}}}[x_3^{2\beta}]\xi_{\id} = [x_3^{2(\alpha+\beta)}]\xi_{\id}$. If $0 \le \alpha + \beta \le 2k$, the statement is clear. If $2k+1 \le \alpha + \beta \le 4k$, applying $D_{F^G}$ to $[-x_3^{2p+2k+1} dx_1 \wedge d x_3]$ and $[x_3^{2p+1}dx_1 \wedge dx_2]$ we have
\[
[2x_2x_3^{2p+2k+1} + x_2^{2(\alpha+\beta)}]_{F^G} = 0,
\]
\[
[2z (p+\frac{1}{2})x_3^{2p} + (2k+1)x_2x_3^{2p+2k+1} + \sum_{\gamma=0}^{2k}2\gamma v^{\gamma}_{\id} x_3^{2(\gamma+p)}]_{F^G} = 0.
\]
These relations give us
\[
[x_3^{2(\alpha+\beta)}]\xi_{\id} = \frac{4}{2k+1}(z (p+\frac{1}{2})x_3^{2p} + \sum_{\gamma=0}^{2k}\gamma v^{\gamma}_{\id} x_3^{2(\gamma+p)}]\xi_{\id}.
\]
This completes the proof.

(ii) $\nabla_{\frac{\partial}{\partial v^0_{g}}}^{(f,G)}[x_1^{2\beta}\xi_{\id}] = \Psi_{\bar{S}}[y_3^{\beta}(y_2+\frac{1}{2}y_3^k)] _{\bar{F}}$. Applying $D_{\bar{F}}$ to $[-y_3^{\beta-1}dy_1\wedge dy_3]$ we obtain:

\[
[2y_3^\beta(y_2+\frac{1}{2}y_3^k)+ \bar{v}_0y_3^{\beta-1}] _{\bar{F}}= 0
\]
from where we calculate
\[
\nabla_{\frac{\partial}{\partial v^0_{g}}}^{(f,G)}[x_1^{2\beta}\xi_{\id}] = \Psi_{\bar{S}}[y_3^{\beta}(y_2+\frac{1}{2}y_3^k)] _{\bar{F}} = \Psi_{\bar{S}}[-\frac{\bar{v}_0}{2}y_3^{\beta-1}] _{\bar{F}}= - \frac{v_g^0}{2}[x_3^{2\beta-2}]\xi_{\id} =  - \frac{v_g^0}{2}\left[\frac{\partial F^G}{\partial v^{\beta-1}_{\id}}\right]\xi_{\id}.
\]
Moreover
\[ 
\nabla_{\frac{\partial}{\partial v^\alpha_{\id}}}^{(f,G)}[\xi_g] = \Psi_{\bar{S}}[y_3^{\beta}(y_2+\frac{1}{2}y_3^k)] _{\bar{F}} = 
\nabla_{\frac{\partial}{\partial v^0_{g}}}^{(f,G)}[x_3^{2\alpha}\xi_{\id}] 
\]
what proves (ii).


(iii) $\nabla_{\frac{\partial}{\partial v^0_{g}}}^{(f,G)}[\xi_{g}] = \Psi_{\bar{S}}[(y_2+\frac{1}{2}y_3^k)^2]_{\bar{F}} = \Psi_{\bar{S}}[y_2^2+ y_2y_3^k + \frac{y_3^{2k}}{4}]_{\bar{F}}$. 
Applying $D_{\bar{F}}$ to $[dy_1\wedge dy_2]$ we obtain:
\[
[y_2^2 + y_2y_3^k + ky_2y_3^k+ \sum_{\gamma = 0}^{2k} \gamma \bar{s}_{\gamma} y_3^{\gamma-1} + \bar v_0\frac{k}{2}y_3^{k-1}]_{\bar{F}} = 0.
\]
Applying $D_{\bar{F}}$ to $[-y_3^{k-1}dy_1\wedge dy_3]$ we obtain:
\[
    [2y_2y_3^k + y_3^{2k} + \bar v_0y_3^{k-1}]_{\bar{F}} = 0.
\]
This gives
\begin{align*}
 [(y_2+\frac{1}{2}y_3^k)^2]_{\bar{F}} &= [-ky_2y_3^k - \sum_{\gamma = 0}^{2k} \gamma \bar{s}_{\gamma} y_3^{\gamma-1} - \bar v_0\frac{k}{2}y_3^{k-1} + \frac{y_3^{2k}}{4}]_{\bar{F}}
 = [ \frac{2k+1}{4} y_3^{2k} - \sum_{\gamma = 0}^{2k} \gamma \bar{s}_{\gamma} y_3^{\gamma-1}]_{\bar{F}}, 
 \\
\nabla_{\frac{\partial}{\partial v^0_{g}}}^{(f,G)}[\xi_{g}] 
    &= \frac{1}{4}[(2k+1)x_3^{4k} - \sum_{\gamma=0}^{2k} 4\gamma v_{\gamma}^{\id} x_3^{2\gamma-2}]\xi_{\id}.
\end{align*}

Note that 
\begin{align*}
 c_gH^F_g = -\frac{H^F_g}{|G|^2} &= -\frac{1}{2k+1} \frac{1}{4} \left( \sum_{\gamma=0}^{2k} v^{\gamma}_{\id} 4\gamma(2\gamma-1)x_3^{2\gamma-1} - (2k+1)^2 x_3^{4k} \right)
 \\
 &=  \frac{1}{4} \left( (2k+1) x_3^{4k} - \sum_{\gamma=0}^{2k} v^{\gamma}_{\id} 4\gamma \frac{2\gamma-1}{2k+1}x_3^{2\gamma-1} \right)
\end{align*}
what completes the proof.

\end{proof}

\end{document}
